% GENERATED FILE. Edit canonical lessons, then run npm run generate:books.
% canonical-source-hash: fnv1a64:5358991492488008
% canonical-lessons: ML-C217-parutti, ML-C217-pattu, ML-C217-tukal, ML-C217-piccala, ML-C217-cunnampu

\chapter{Cotton, Silk, Leather, Brass, Lime}
\label{ch:ml-cotton-silk-leather-brass-lime}
\hlchaptermodality{217}

\begin{chapteropening}
\textbf{By the end of this chapter:} \emph{I can say cotton, silk, leather, brass and lime.}

Complete the last lesson of chapter 217: I can say cotton, silk, leather, brass and lime.
\end{chapteropening}

\section[\texorpdfstring{\ml{പരുത്തി} (parutti)}{പരുത്തി (parutti)}]{\ml{പരുത്തി} (parutti) — cotton}

\label{lesson:ML-C217-parutti}



\begin{quote}
Before the new one: say the Malayalam for a brick, then the Malayalam for cement.
\end{quote}

\subsection*{You'll want to know: \ml{പരുത്തി}}
\textbf{\ml{പരുത്തി}} — \emph{parutti} — \textquotedblleft{}cotton\textquotedblright{}.

\begin{grammarlens}[title={What you've built}]
One more word for this chapter.
\end{grammarlens}

\subsection*{Your turn}
\begin{itemize}
\raggedright
  \item \emph{Say it:} \emph{parutti}
  \item \emph{Say it:} \emph{parutti}, once more
  \item \emph{Say it:} say \emph{parutti}
  \item \emph{Recall:} say the Malayalam for a brick, then the Malayalam for cement, then say \emph{parutti} again
\end{itemize}

\begin{tcolorbox}[breakable,colback=teal!4,colframe=teal!35!black,title={Before you move on}]
What does \ml{പരുത്തി} mean? (\textquotedblleft{}Cotton\textquotedblright{}.) Say it once more.
\end{tcolorbox}
\section[\texorpdfstring{\ml{പട്ട്} (paṭṭŭ)}{പട്ട് (paṭṭŭ)}]{\ml{പട്ട്} (paṭṭŭ) — silk}

\label{lesson:ML-C217-pattu}



\begin{quote}
Before the new one: say the Malayalam for cement, then the Malayalam for cotton.
\end{quote}

\subsection*{You'll want to know: \ml{പട്ട്}}
\textbf{\ml{പട്ട്}} — \emph{paṭṭŭ} — \textquotedblleft{}silk\textquotedblright{}.

\begin{grammarlens}[title={What you've built}]
One more word for this chapter.
\end{grammarlens}

\subsection*{Your turn}
\begin{itemize}
\raggedright
  \item \emph{Say it:} \emph{paṭṭŭ}
  \item \emph{Say it:} \emph{paṭṭŭ}, once more
  \item \emph{Say it:} say \emph{paṭṭŭ}
  \item \emph{Recall:} say the Malayalam for cement, then the Malayalam for cotton, then say \emph{paṭṭŭ} again
\end{itemize}

\begin{tcolorbox}[breakable,colback=teal!4,colframe=teal!35!black,title={Before you move on}]
What does \ml{പട്ട്} mean? (\textquotedblleft{}Silk\textquotedblright{}.) Say it once more.
\end{tcolorbox}
\section[\texorpdfstring{\ml{തുകൽ} (tukal)}{തുകൽ (tukal)}]{\ml{തുകൽ} (tukal) — leather}

\label{lesson:ML-C217-tukal}



\begin{quote}
Before the new one: say the Malayalam for cotton, then the Malayalam for silk.
\end{quote}

\subsection*{You'll want to know: \ml{തുകൽ}}
\textbf{\ml{തുകൽ}} — \emph{tukal} — \textquotedblleft{}leather\textquotedblright{}.

\begin{grammarlens}[title={What you've built}]
One more word for this chapter.
\end{grammarlens}

\subsection*{Your turn}
\begin{itemize}
\raggedright
  \item \emph{Say it:} \emph{tukal}
  \item \emph{Say it:} \emph{tukal}, once more
  \item \emph{Say it:} say \emph{tukal}
  \item \emph{Recall:} say the Malayalam for cotton, then the Malayalam for silk, then say \emph{tukal} again
\end{itemize}

\begin{tcolorbox}[breakable,colback=teal!4,colframe=teal!35!black,title={Before you move on}]
What does \ml{തുകൽ} mean? (\textquotedblleft{}Leather\textquotedblright{}.) Say it once more.
\end{tcolorbox}
\section[\texorpdfstring{\ml{പിച്ചള} (piccaḷa)}{പിച്ചള (piccaḷa)}]{\ml{പിച്ചള} (piccaḷa) — brass}

\label{lesson:ML-C217-piccala}



\begin{quote}
Before the new one: say the Malayalam for silk, then the Malayalam for leather.
\end{quote}

\subsection*{You'll want to know: \ml{പിച്ചള}}
\textbf{\ml{പിച്ചള}} — \emph{piccaḷa} — \textquotedblleft{}brass\textquotedblright{}.

\begin{grammarlens}[title={What you've built}]
One more word for this chapter.
\end{grammarlens}

\subsection*{Your turn}
\begin{itemize}
\raggedright
  \item \emph{Say it:} \emph{piccaḷa}
  \item \emph{Say it:} \emph{piccaḷa}, once more
  \item \emph{Say it:} say \emph{piccaḷa}
  \item \emph{Recall:} say the Malayalam for silk, then the Malayalam for leather, then say \emph{piccaḷa} again
\end{itemize}

\begin{tcolorbox}[breakable,colback=teal!4,colframe=teal!35!black,title={Before you move on}]
What does \ml{പിച്ചള} mean? (\textquotedblleft{}Brass\textquotedblright{}.) Say it once more.
\end{tcolorbox}
\section[\texorpdfstring{\ml{ചുണ്ണാമ്പ്} (cuṇṇāmpŭ)}{ചുണ്ണാമ്പ് (cuṇṇāmpŭ)}]{\ml{ചുണ്ണാമ്പ്} (cuṇṇāmpŭ) — lime (for whitewash)}

\label{lesson:ML-C217-cunnampu}



\begin{quote}
Before the new one: say the Malayalam for leather, then the Malayalam for brass.
\end{quote}

\subsection*{You'll want to know: \ml{ചുണ്ണാമ്പ്}}
\textbf{\ml{ചുണ്ണാമ്പ്}} — \emph{cuṇṇāmpŭ} — \textquotedblleft{}lime (for whitewash)\textquotedblright{}.

\begin{grammarlens}[title={What you've built}]
One more word for this chapter.
\end{grammarlens}

\subsection*{Your turn}
\begin{itemize}
\raggedright
  \item \emph{Say it:} \emph{cuṇṇāmpŭ}
  \item \emph{Say it:} \emph{cuṇṇāmpŭ}, once more
  \item \emph{Say it:} say \emph{cuṇṇāmpŭ}
  \item \emph{Recall:} say the Malayalam for leather, then the Malayalam for brass, then say \emph{cuṇṇāmpŭ} again
\end{itemize}

\begin{tcolorbox}[breakable,colback=teal!4,colframe=teal!35!black,title={Before you move on}]
What does \ml{ചുണ്ണാമ്പ്} mean? (\textquotedblleft{}Lime (for whitewash)\textquotedblright{}.) Say it once more.
\end{tcolorbox}
